\section{Challenges and Solutions}\label{sec:challenges}
% The brief asks for a one-page reflection on the network challenges faced and how they were solved.

The hardest problems were not syntax errors but designs that are correct in normal operation and fail
only under a particular failure. We found them by running the routing logic before the lab
(\cref{sec:emulation}) and by reasoning about each failure scenario. Each fix is now guarded by an
automated test, so a later change cannot reintroduce it.

\textbf{1. A routing loop through the BGP-free core.} \emph{Symptom:} when Hanoi lost its ISP-A connection,
traffic from HQ to Hanoi looped until its TTL expired. \emph{Cause:} \dev{HCMC-EDGE-RTR01} still had Internet
access, so it kept attracting all traffic with its default route; for Hanoi it now had to forward to
\dev{HCMC-EDGE-RTR02}. The only path between the two edges ran through a core switch, which knows no branch
prefixes and returned the packet on its default route. \emph{Fix:} a direct routed link between the two HQ
edge routers, so every BGP next hop is reachable without crossing the core. \emph{Evidence:} emulation
scenario S2 passes with the link; S9 reproduces the loop without it.

\textbf{2. A dead router that still looked reachable.} \emph{Symptom:} after \dev{HCMC-EDGE-RTR01} failed,
\dev{HCMC-EDGE-RTR02} kept using the paths it had learned from it and discarded the traffic. \emph{Cause:} BGP
removes a route when its next hop becomes unreachable, but the dead router's loopback lies inside the site
summary \texttt{10.1.0.0/16}, which points to \texttt{Null0}; the next hop kept ``resolving'' and only the
default hold time of 180\,s would have removed the routes. \emph{Fix:} 3\,s/9\,s timers on that iBGP session
as well. \emph{Lesson:} a summary route to \texttt{Null0} defeats next-hop tracking for every address it
covers.

\textbf{3. Failover that waited.} The standby gateway had the same 60\,s preempt delay as the active one,
so a tracking-driven takeover was postponed by a minute, and IOS polls route-tracking objects only every 15\,s.
We kept the delay on the designed owner only and set the polling interval to 1\,s. In the same review we
lowered the OSPF dead interval from 40\,s to 4\,s, because virtual links -- like any link through an
intermediate switch -- do not go down when the neighbour dies.

\textbf{4. IPv6 that could not cross the overlay.} The first version used IPsec virtual tunnel interfaces in
\texttt{ipsec ipv4} mode, which carry IPv4 only, while the design runs IPv6 BGP sessions across the tunnels.
GRE inside IPsec carries both address families in one tunnel. Two related corrections followed: each IPv6
gateway got its own HSRP group, and the IPv6 IoT access list now permits neighbour discovery before its final
deny.

\textbf{5. An address plan that overlapped itself.} The first reserved growth blocks overlapped the
infrastructure ranges for loopbacks and links. The address-plan test found it on the first run, before a
single device was configured; the reserved blocks were reduced to aligned blocks outside those ranges.

\todo{add the problems met while building the lab (e.g.\ commands rejected by the images, memory limits,
console paste errors) and how you solved them -- keep the whole section to one page}
